% page 372 of the facsimile (image p-371 of the scan)
% block 154.9 x 229.0 mm ; left margin 8.0 mm
\pgc{-12.700mm}{0.000mm}{
\l{\cel{3}364}
\l{}
\l{\sou{melanismo}. (\sou{patol.}) Stando di la pelo qua, loke o tota-}
\l{\cel{3}korpe, aquiras koloro brunatra o nigratra, kun intens-}
\l{\cel{3}eso-grado variiva, ma (omnakaze) ne precize inkluzita.}
\l{}
\l{\sou{melaso.}\cel{2}Reziduo siropatra de la kristalesko e rafineso}
\l{\cel{3}di sukro. - DEFIS.}
\l{}
\l{\sou{meleo.}\cel{1}(en\cel{1}\sou{kombato}).Intermixeso sen-ordina, pelmele, di la}
\l{\cel{3}kombatanti qui frapas reciproke. - EF.}
\l{}
\l{\sou{melinito}. (\sou{kemio)}. Mixuro exploziva, di qua la bazo esas tri-}
\l{\cel{3}nitrofenolo 2.4.6.\cel{2}- DEFIRS.}
\l{}
\l{\sou{meliso}. (\sou{bot.}) Planto aromata, ek la familio \textquotedbl{}labiozi\textquotedbl{}, di}
\l{\cel{3}qua la folii uzesas por la faco di aquo alkoholizita e}
\l{\cel{3}stomako-medikamenta. - DEFIRS.}
\l{}
\l{\sou{melkar}. (\sou{trans.}) Extraktar la lakto di femino, per presar}
\l{\cel{3}lua mamo. - DE.}
\l{}
\l{\sou{melodio.}I.(\sou{muziko}).Intersequo de soni qui plezas a la au-}
\l{\cel{3}danti od askoltanti. - II. Intersequo de vorti, de frazi}
\l{\cel{3}qui plezas a la audanti od askoltanti. - DEFIRS.}
\l{}
\l{\sou{melodramo.}\cel{2}I. Dramato, akompanata da muziko. - II. Dra-}
\l{\cel{3}mato populara, qua probas emocigar per la extremeso}
\l{\cel{3}di la cirkonstanci e la exajero di la sentimenti.DEFIRS}
\l{}
\l{\sou{melolonto.(zool.)}\cel{2}Insekto koleoptera, qua aparas cirkum}
\l{\cel{3}la fino di aprilo, e di qua la larvo devoras la radiki}
\l{\cel{3}di la planti. - L.\cel{1}\sou{melolontha vulgaris}. - IL.}
\l{}
\l{\sou{melono.}\cel{2}(\sou{bot.)}\cel{2}Planto ek la familio \textquotedbl{}kukurbitacei\textquotedbl{},di}
\l{\cel{3}qua la frukto esas sukoza e sukroza. - L.\cel{1}\sou{cucumis melo}.}
\l{\cel{3}- DEFS.}
\l{}
\l{\sou{melopear}. (\sou{netrans.}) Kantar ritmoze e segun-mezure, por}
\l{\cel{3}akompanar la deklamo parolata. - EFIRS.}
\l{}
\l{\sou{mem}. (\sou{gram.}) -. Manier-adverbo uzata (1) avan komparativo}
\l{\cel{3}(di supereso o di infreso) por plufortigar lu ; (2)}
\l{\cel{3}avan substantivo o verbo, por insistar. - F.}
\l{}
\l{\sou{membro}. I. Apendico unionita a la torso di la homo, di la}
\l{\cel{3}animalo, per artiki. - II. Persono qua esas un ek la}
\l{\cel{3}kompozanti di korporaciono politikala, literaturala -}
\l{\cel{3}di asociitaro, di komuniono religiala. - III. Un ek la}
\l{\cel{3}propozicioni ye qui konsistas frazo ; un ek la frazi}
\l{\cel{3}ye qui konsistas la ideo-developo. - IV. Singla de la}
\l{\cel{3}du parti di equaciono, quin separas la signo di inter-}
\l{\cel{3}egaleso. - EFIS.}
}
